\documentclass[10pt, letterpaper]{article}

% Packages:
\usepackage[
    ignoreheadfoot, % set margins without considering header and footer
    top=2 cm, % seperation between body and page edge from the top
    bottom=2 cm, % seperation between body and page edge from the bottom
    left=2 cm, % seperation between body and page edge from the left
    right=2 cm, % seperation between body and page edge from the right
    footskip=1.0 cm, % seperation between body and footer
    % showframe % for debugging 
]{geometry} % for adjusting page geometry
\usepackage{titlesec} % for customizing section titles
\usepackage{tabularx} % for making tables with fixed width columns
\usepackage{array} % tabularx requires this
\usepackage[dvipsnames]{xcolor} % for coloring text
\definecolor{primaryColor}{RGB}{0, 0, 0} % define primary color
\usepackage{enumitem} % for customizing lists
\usepackage{fontawesome5} % for using icons
\usepackage{amsmath} % for math
\usepackage[
    pdftitle={Yichao Yuan's CV},
    pdfauthor={Yichao Yuan},
    pdfcreator={LaTeX with RenderCV},
    colorlinks=true,
    urlcolor=primaryColor
]{hyperref} % for links, metadata and bookmarks
\usepackage[pscoord]{eso-pic} % for floating text on the page
\usepackage{calc} % for calculating lengths
\usepackage{bookmark} % for bookmarks
\usepackage{lastpage} % for getting the total number of pages
\usepackage{changepage} % for one column entries (adjustwidth environment)
\usepackage{paracol} % for two and three column entries
\usepackage{ifthen} % for conditional statements
\usepackage{needspace} % for avoiding page brake right after the section title
\usepackage{iftex} % check if engine is pdflatex, xetex or luatex

% Ensure that generate pdf is machine readable/ATS parsable:
\ifPDFTeX
    \input{glyphtounicode}
    \pdfgentounicode=1
    \usepackage[T1]{fontenc}
    \usepackage[utf8]{inputenc}
    \usepackage{lmodern}
\fi

\usepackage{charter}

% Some settings:
\raggedright
\AtBeginEnvironment{adjustwidth}{\partopsep0pt} % remove space before adjustwidth environment
\pagestyle{empty} % no header or footer
\setcounter{secnumdepth}{0} % no section numbering
\setlength{\parindent}{0pt} % no indentation
\setlength{\topskip}{0pt} % no top skip
\setlength{\columnsep}{0.15cm} % set column seperation
\pagenumbering{gobble} % no page numbering

\titleformat{\section}{\needspace{4\baselineskip}\bfseries\large}{}{0pt}{}[\vspace{1pt}\titlerule]

\titlespacing{\section}{
    % left space:
    -1pt
}{
    % top space:
    0.3 cm
}{
    % bottom space:
    0.2 cm
} % section title spacing

\renewcommand\labelitemi{$\vcenter{\hbox{\small$\bullet$}}$} % custom bullet points
\newenvironment{highlights}{
    \begin{itemize}[
        topsep=0.10 cm,
        parsep=0.10 cm,
        partopsep=0pt,
        itemsep=0pt,
        leftmargin=0 cm + 10pt
    ]
}{
    \end{itemize}
} % new environment for highlights


\newenvironment{highlightsforbulletentries}{
    \begin{itemize}[
        topsep=0.10 cm,
        parsep=0.10 cm,
        partopsep=0pt,
        itemsep=0pt,
        leftmargin=10pt
    ]
}{
    \end{itemize}
} % new environment for highlights for bullet entries

\newenvironment{onecolentry}{
    \begin{adjustwidth}{
        0 cm + 0.00001 cm
    }{
        0 cm + 0.00001 cm
    }
}{
    \end{adjustwidth}
} % new environment for one column entries

\newenvironment{twocolentry}[2][]{
    \onecolentry
    \def\secondColumn{#2}
    \setcolumnwidth{\fill, 4.5 cm}
    \begin{paracol}{2}
}{
    \switchcolumn \raggedleft \secondColumn
    \end{paracol}
    \endonecolentry
} % new environment for two column entries

\newenvironment{threecolentry}[3][]{
    \onecolentry
    \def\thirdColumn{#3}
    \setcolumnwidth{, \fill, 4.5 cm}
    \begin{paracol}{3}
    {\raggedright #2} \switchcolumn
}{
    \switchcolumn \raggedleft \thirdColumn
    \end{paracol}
    \endonecolentry
} % new environment for three column entries

\newenvironment{header}{
    \setlength{\topsep}{0pt}\par\kern\topsep\centering\linespread{1.5}
}{
    \par\kern\topsep
} % new environment for the header

\newcommand{\placelastupdatedtext}{% \placetextbox{<horizontal pos>}{<vertical pos>}{<stuff>}
  \AddToShipoutPictureFG*{% Add <stuff> to current page foreground
    \put(
        \LenToUnit{\paperwidth-2 cm-0 cm+0.05cm},
        \LenToUnit{\paperheight-1.0 cm}
    ){\vtop{{\null}\makebox[0pt][c]{
        \small\color{gray}\textit{Last updated in September 2024}\hspace{\widthof{Last updated in September 2024}}
    }}}%
  }%
}%

% save the original href command in a new command:
\let\hrefWithoutArrow\href

% new command for external links:


\begin{document}
    \newcommand{\AND}{\unskip
        \cleaders\copy\ANDbox\hskip\wd\ANDbox
        \ignorespaces
    }
    \newsavebox\ANDbox
    \sbox\ANDbox{$|$}

    \begin{header}
        \fontsize{25 pt}{25 pt}\selectfont Yichao Yuan

        \vspace{5 pt}

        \normalsize
        \mbox{Champaign, IL}%
        \kern 5.0 pt%
        \AND%
        \kern 5.0 pt%
        \mbox{\hrefWithoutArrow{mailto:yichaoy2@illinois.edu}{yichaoy2@illinois.edu}}%
        \kern 5.0 pt%
        \AND%
        \kern 5.0 pt%
        \mbox{\hrefWithoutArrow{tel:1-734-450-4291}{734 450 4291}}%
        \kern 5.0 pt%
        \AND%
        \kern 5.0 pt%
        \mbox{\hrefWithoutArrow{https://yichao-yuan-99.github.io/}{yichao-yuan-99.github.io}}%
        \kern 5.0 pt%
        \AND%
        \kern 5.0 pt%
        \mbox{\hrefWithoutArrow{https://www.linkedin.com/in/yichao-yuan-30946b220/}{linkedin.com/in/yichao-yuan-30946b220}}%
        \kern 5.0 pt%
        \AND%
        \kern 5.0 pt%
        \mbox{\hrefWithoutArrow{https://github.com/yichao-yuan-99}{github.com/yichao-yuan-99}}%
    \end{header}

    \vspace{5 pt - 0.3 cm}


    % \section{Welcome to RenderCV!}



        
    %     \begin{onecolentry}
    %         \href{https://rendercv.com}{RenderCV} is a LaTeX-based CV/resume version-control and maintenance app. It allows you to create a high-quality CV or resume as a PDF file from a YAML file, with \textbf{Markdown syntax support} and \textbf{complete control over the LaTeX code}.
    %     \end{onecolentry}

    %     \vspace{0.2 cm}

    %     \begin{onecolentry}
    %         The boilerplate content was inspired by \href{https://github.com/dnl-blkv/mcdowell-cv}{Gayle McDowell}.
    %     \end{onecolentry}


    
    % \section{Quick Guide}

    % \begin{onecolentry}
    %     \begin{highlightsforbulletentries}


    %     \item Each section title is arbitrary and each section contains a list of entries.

    %     \item There are 7 unique entry types: \textit{BulletEntry}, \textit{TextEntry}, \textit{EducationEntry}, \textit{ExperienceEntry}, \textit{NormalEntry}, \textit{PublicationEntry}, and \textit{OneLineEntry}.

    %     \item Select a section title, pick an entry type, and start writing your section!

    %     \item \href{https://docs.rendercv.com/user_guide/}{Here}, you can find a comprehensive user guide for RenderCV.


    %     \end{highlightsforbulletentries}
    % \end{onecolentry}

    \section{Education}
        \begin{onecolentry}
            \textbf{University of Illinois Urbana-Champaign}, Champaign, Illinois
        \end{onecolentry}
        \vspace{0.10cm}
        \begin{twocolentry} {
            Aug 2025 - May 2027 (Expected)
        }
        Ph.D. Student in Computer Science
        \vspace{0.10cm}

        \textit{Area of focus: GPU-accelerated system design, System/Architecture design for AI applications, System/Architecture-Algorithm Co-Design}
        
        \end{twocolentry}
        
        \begin{onecolentry}
            \textbf{University of Michigan}, Ann Arbor, Michigan
        \end{onecolentry}
        \vspace{0.10cm}
        \begin{twocolentry} {
            Aug 2023 - Aug 2025
        }
        Ph.D. Student in Computer Science and Engineering, 4.0/4.0

        \textit{Transferred to University of Illinois Urbana-Champaign with advisor (Aug 2025)}
        \end{twocolentry}
        \begin{twocolentry} {
            Aug 2021 - Apr 2023
        }
        M.S.E. in Electrical and Computer Engineering, 4.0/4.0        \end{twocolentry}

        
        \vspace{0.15cm}
        \begin{onecolentry}
            \textbf{Shanghai Jiao Tong University (SJTU)}, Shanghai, China\end{onecolentry}
        \begin{onecolentry}
            University of Michigan -- Shanghai Jiao Tong University Joint Institute (UM-SJTU JI)
        \end{onecolentry}
        \vspace{0.10cm}
        
        \begin{twocolentry}{
            Sept 2017 - Aug 2021
        }
            B.S.E. in Electrical and Computer Engineering, 3.8/4.0\end{twocolentry}
        \vspace{0.10cm}



    \section{Publications}
        \textbf{Y. Yuan}, M. Chowdhury, and N. Talati, “KAIROS: Stateful, Context-Aware Power-Efficient Agentic Inference Serving,” \textbf{ASPLOS 2027}.

        \vspace{0.25cm}
        
        \textbf{Y. Yuan}, A. Nayak, S. Kundu, and N. Talati, “Agentic AI Workload Characteristics,” \textbf{IISWC 2026}.
    
        \vspace{0.25cm}
    

        \textbf{Y. Yuan}, L. Ma, and N. Talati, “MoE-Lens: Towards the Hardware Limit of High-Throughput MoE LLM Serving Under Resource Constraints,” at the 34th International Symposium on High-Performance Parallel and Distributed Computing, \textbf{HPDC 2026}.

        \vspace{0.25cm}

        Y. Xia, D. Sharma, \textbf{Y. Yuan}, S. Kundu, and N. Talati, “MoDM: Efficient Serving for Image Generation via Mixture-of-Diffusion Models,”  at 31st Conference on Architectural Support for Programming Languages and Operating Systems \textbf{ASPLOS 2026}.
        
        \vspace{0.25cm}

        \textbf{Y. Yuan}, A. Iyer, L. Ma, N. Talati, ``Vortex: Overcoming Memory Capacity Limitations in GPU-Accelerated Large-Scale Data Analytics,'' at International Conference on Very Large Databases \textbf{VLDB 2025}

        \vspace{0.25cm}

        H. Ye, Y. Xia, Y. Chen, K-Y. Chen, \textbf{Y. Yuan}, S. Deng, B. Kasikci, T. Mudge, and N. Talati, ``Palermo: Improving the Performance of Oblivious Memory using Protocol-Hardware Co-Design,'' at the 31st IEEE International Symposium on High-Performance Computer Architecture \textbf{HPCA 2025}.
        
        \vspace{0.25cm}

        J. Pavon, I. Valdivieso, C. Morales, C. Hernandez, M. Aslan, J. Lindegger, \textbf{Y. Yuan}, R. Bagué, M. Alser, O. Mutlu, S. Marco-Sola, O. Ergin, N. Talati, M. Valero, O. Unsal, A. Cristal, ``QUETZL: Vector Acceleration Framework for Modern Genome Sequence Analysis Algorithms,'' at 2024 International Symposium on Computer Architecture \textbf{ISCA 2024}.

        \vspace{0.25cm}

        \textbf{Y. Yuan}, H. Ye, S. Vedula, W. Kaza, and N. Talati, "Everest: GPU-Accelerated System for Mining Temporal Motifs," at International Conference on Very Large Databases \textbf{VLDB 2024}.

        \vspace{0.25cm}

         N. Talati, H. Ye, S. Vedula, K-Y Chen, Y. Chen, D. Liu, \textbf{Y. Yuan}, D. Blaauw, A. Bronstein, T. Mudge, and R. Dreslinski, "Mint: An Accelerator For Mining Temporal Motifs," at the 55th IEEE/ACM International Symposium on Microarchitecture \textbf{MICRO 2022}.

    \section{Professional Experience}
    
        \begin{twocolentry}{
        May 2026 - Aug 2026
        }
        \textbf{Software Engineering Intern}, Google -- San Jose, CA
        \end{twocolentry}
        
        \vspace{0.10 cm}

        \begin{onecolentry}
        \begin{highlights}
        \item Investigated GPU-accelerated query processing, evaluated how kernel interface design affects GPU operator performance through microarchitectural resource utilization and interface-induced overheads, and developed and tested techniques for query-level acceleration.
        \end{highlights}
        \end{onecolentry}
        
        \vspace{0.2 cm}
        
        \begin{twocolentry}{
        May 2025 - Aug 2025
        }
        \textbf{System Architecture Intern}, Arm -- Austin, TX
        \end{twocolentry}
        
        \vspace{0.10 cm}
        
        \begin{onecolentry}
        \begin{highlights}
        \item Researched memory-access performance in chiplet architectures and investigated architectural support and kernel performance for KV-cache management and attention operations in CPU-based LLM inference.
        \end{highlights}
        \end{onecolentry}
        
        \vspace{0.2 cm}




    \section{Professional Service}

    Review Board Member, VLDB 2027
    
    Program Committee Member, VLDB 2026 Demonstration Track
    
    Shadow Program Committee Member, VLDB 2026
    
    Artifact Evaluation Committee Member, ASPLOS 2025
    
    Artifact Evaluation Committee Member, ISCA 2025

    
    \section{Technical Skills}


        \begin{onecolentry}
            \textbf{Languages:} C++, C, Python, CUDA, HIP, SystemVerilog
        \end{onecolentry}

        \vspace{0.2 cm}

        \begin{onecolentry}
            \textbf{Tools/frameworks:} NVIDIA Nsight Systems/Compute, ROCm profiler, PyTorch, OpenMP, x86/ARM vector intrinsics
        \end{onecolentry}


    

\end{document}